\icode{E-MC-SI-L3-261026-002}
一质量为\(m\)的物体在水平面上运动，它的两个正交分速度\(v_x\)、\(v_y\)随时间变化的图线分别如图甲、乙所示，由图可知\paren[]

\begin{tikzpicture}[>=stealth,font=\scriptsize]
	\draw[->,blue] (0,0)node[below left]{ \(O\)}--(3.5,0)node[below]{ \(t\)/\qty{}{s}};
	\draw[->,blue] (0,0)--(0,2.5)node[right]{ \(v_x\)/(\qty{}{m\cdot s^{-1}})};
	\foreach \x/\n in {.5/1,1/2,1.5/3,2/4,2.5/5,3/6}
		\draw[blue] (\x,.1)--(\x,0)node[below]{ \(\n\)};
	\foreach \y/\n in {.5/1,1/2,1.5/3,2/4}
		\draw[blue] (.1,\y)--(0,\y)node[left]{ \(\n\)};
	\draw[dashed,cyan] (0,2)--(2,2)--(2,0);
	\draw[magenta] (0,0)--(2,2)--(3,0);
	\node at(1.75,-.7){ 图甲};
\end{tikzpicture}
\begin{tikzpicture}[>=stealth,font=\scriptsize]
	\draw[->,blue] (0,0)node[below left]{ \(O\)}--(3.5,0)node[below]{ \(t\)/\qty{}{s}};
	\draw[->,blue] (0,0)--(0,2.5)node[right]{ \(v_y\)/(\qty{}{m\cdot s^{-1}})};
	\foreach \x/\n in {.5/1,1/2,1.5/3,2/4,2.5/5,3/6}
		\draw[blue] (\x,.05)--(\x,0)node[below]{ \(\n\)};
	\foreach \y/\n in {.5/1,1/2,1.5/3,2/4}
		\draw[blue] (.05,\y)--(0,\y)node[left]{ \(\n\)};
	\draw[dashed,cyan] (2,1)--(2,0);
	\draw[magenta] (0,1)--(2,1)--(3,0);
	\node at(1.75,-.7){ 图乙};
\end{tikzpicture}
\begin{choices}
	\item 从开始至\qty{6}{s}末一直做直线运动
	\item 从开始至\qty{6}{s}末一直做曲线运动
	\item \qty{4}{s}末物体的速度大小为\(2\sqrt{5}\) \qty{}{m/s}
	\item 最初\qty{4}{s}的位移为\qty{16}{m}
\end{choices}
\begin{answer}
	C\qquad 提示：容易误选B，\(4\sim\)\qty{6}{s}是匀减速直线运动
\end{answer}